% Рисунок с доски: четырёхмерное пространство-время. Характеристический конус
% с вершиной (x,t); срез на уровне tau -- сфера S_{x,a(t-tau)}.
\begin{figure}[htbp]
\centering
\begin{tikzpicture}[>={Latex[length=4pt]}]
  \draw[->] (0,-0.5) -- (0,3.9) node[anchor=south east] {$t$};
  \draw[->] (-0.5,0) -- (6.3,0) node[anchor=north east] {$x \in \R^3$};

  % боковая поверхность конуса
  \draw[thick, blue!60!black] (3,3.0) -- (1,0);
  \draw[thick, blue!60!black] (3,3.0) -- (5,0);

  % основание: сфера S_{x,at}
  \draw[thick, blue!60!black] (3,0) ellipse [x radius=2, y radius=0.34];

  % срез на уровне tau
  \draw[thick, red!65!black] (3,1.3) ellipse [x radius=1.133, y radius=0.2];

  \fill (3,3.0) circle (1.6pt);
  \node[anchor=south, font=\footnotesize] at (3,3.08) {$(x,t)$};
  \fill (3,0) circle (1.4pt);
  \node[anchor=north east, font=\footnotesize] at (3.02,-0.06) {$x$};

  \draw[dashed, gray] (0,3.0) -- (3,3.0);
  \draw[dashed, gray] (0,1.3) -- (1.867,1.3);
  \node[anchor=east, font=\footnotesize] at (-0.06,3.0) {$t$};
  \node[anchor=east, font=\footnotesize] at (-0.06,1.3) {$\tau$};

  \draw[->, red!65!black] (5.5,1.75) -- (4.18,1.38);
  \node[red!65!black, anchor=west, font=\footnotesize] at (5.55,1.78)
    {$S_{x,a(t-\tau)}$};
  \draw[->, blue!60!black] (5.9,0.95) -- (4.35,0.6);
  \node[blue!60!black, anchor=west, font=\footnotesize] at (5.95,0.98)
    {область зависимости};
  \draw[->, blue!60!black] (1.1,-0.95) -- (2.1,-0.28);
  \node[blue!60!black, anchor=north, font=\footnotesize] at (1.1,-1.0)
    {$S_{x,at}$};

  \node[anchor=west, font=\footnotesize] at (4.6,3.3)
    {4-мерное пространство-время};
\end{tikzpicture}
\caption{Характеристический конус с вершиной $(x,t)$. Значение $u(x,t)$
определяется данными не во всём конусе, а только на его \emph{боковой
поверхности}: начальными данными на сфере $S_{x,at}$ и правой частью на сферах
$S_{x,a(t-\tau)}$, $0 \le \tau \le t$.}
\label{fig:cone-kirchhoff}
\end{figure}
